\documentclass[11pt,a4paper]{article}

\usepackage[utf8]{inputenc}
\usepackage[T1]{fontenc}
\usepackage{geometry}
\usepackage{amsmath}
\usepackage{longtable}
\usepackage{array}
\usepackage{url}

\geometry{margin=2.5cm}

\newcommand{\toprule}{\hline}
\newcommand{\midrule}{\hline}
\newcommand{\bottomrule}{\hline}

\title{Technical Report: AdaptiveRoute\\
\large Agentic Route Replanning with a Fine-Tuned Policy and an Exact Solver}
\author{Gabriel Carvalho Domingos da Conceição}
\date{August 29, 2026}

\begin{document}

\maketitle

\begin{abstract}
This report documents the design, implementation and evaluation of AdaptiveRoute, a system for
replanning last-mile delivery routes when operational disruptions occur mid-execution. The system
converts a driver's natural-language report --- a blocked road, an unavailable customer --- into a
validated route plan. It combines three components with strictly separated responsibilities: a
general-purpose language model that interprets the message and orchestrates the workflow, a
fine-tuned 7B policy that proposes route candidates, and an exact Pyomo + HiGHS optimizer that
remains the sole authority on feasibility. No plan reaches a driver without passing deterministic
validation, regardless of which component produced it. Around this core the system carries the
supporting AI-engineering layer that makes a conversational interface viable: a bounded rolling
context window, a vector-backed retrieval layer over project documentation, and a grounded route
question-answering pipeline. The report covers the architecture, the agentic workflow, the validation
contract, these components, the supervised fine-tuning of the routing policy, and the measured limits
of the resulting system.
\end{abstract}

\tableofcontents

\section{Executive Summary}

AdaptiveRoute addresses a gap that is neither an optimization problem nor a language problem, but
the handoff between them. A delivery route optimized overnight becomes invalid during execution;
re-running the optimizer requires someone to translate a disruption into a formal model, and asking
a language model directly produces plausible routes that violate capacity or skip customers.

The system closes that handoff with a three-layer architecture:

\begin{itemize}
  \item a general-purpose LLM interprets the driver's message and extracts a structured operational
        event;
  \item a fine-tuned LoRA policy proposes a route candidate in approximately 1.7 seconds;
  \item deterministic validation accepts, repairs, or rejects that candidate, with Pyomo + HiGHS as
        the fallback authority.
\end{itemize}

The fine-tuned routing policy, referred to as v5, reaches a 94.4\% feasible-candidate rate on a fixed
1,000-example benchmark. The remaining 5.6\% is not a correctness risk, because the architecture
never depends on the model being right: invalid candidates are repaired or replaced by the solver.
The measured operating boundary of the policy is 8--9 customers per scenario; beyond that the
cascade falls through to the solver on every request.

The selected adapter is
\path{outputs/models/adaptiveroute-qwen2_5-7b-lora-error20k-v5}, based on
\path{Qwen/Qwen2.5-7B-Instruct}.

\section{Project Objective and Scope}

The objective was to build a working end-to-end system, not a demonstration of isolated components.
The scope covers:

\begin{itemize}
  \item exact CVRP planning over ingested orders;
  \item natural-language reporting of operational disruptions;
  \item agentic replanning with validation, repair and fallback;
  \item conversation memory and route question answering;
  \item an operational console for administrators and a scoped workspace for drivers.
\end{itemize}

The routing scenarios are synthetic CVRP instances, generated primarily with 8 customers and 2
vehicles. The target behaviour of the learned component is operational feasibility, not proof of
optimality: a generated plan is useful only if it passes deterministic validation.

\section{Theoretical Foundation}

\subsection{Vehicle Routing and CVRP}

The Vehicle Routing Problem was introduced by Dantzig and Ramser in the classical truck dispatching
formulation \cite{dantzig1959}. In the Capacitated Vehicle Routing Problem, a fleet of vehicles with
limited capacity must serve a set of customers from a depot while minimizing total travel distance.

CVRP is NP-hard. For realistic instances, exact optimization becomes expensive, and practical
systems combine mathematical programming, heuristics and learned policies.

\subsection{Prescriptive Optimization and Subtour Elimination}

The prescriptive layer answers: given a scenario and a disruption, what route plan should be
executed? Pyomo was selected because it allows explicit algebraic modeling in Python and keeps the
optimization model inspectable, testable and version-controlled \cite{pyomo}. HiGHS was selected as
the open-source solver for LP and MIP problems, with a Python interface suitable for local
experimentation \cite{highs}.

Subtour elimination uses the Miller--Tucker--Zemlin formulation \cite{mtz}. MTZ is compact and
readable, requiring only $O(n^2)$ additional constraints, which suits a system where the model is
meant to be auditable. Its known weakness is a loose linear relaxation compared with
Dantzig--Fulkerson--Johnson cut families, which becomes the binding constraint on scalability as
instance size grows. This trade-off is revisited in Section~\ref{sec:limitations}.

\subsection{Learning for Combinatorial Optimization}

A substantial literature asks whether neural networks can solve routing problems directly. Pointer
Networks introduced an attention-based architecture able to output permutations over variable-length
inputs, and demonstrated it on TSP \cite{pointer}. Kool, van Hoof and Welling extended this line with
an attention model trained by reinforcement learning that produces high-quality tours for TSP and
VRP without a solver in the loop \cite{kool2019}.

These approaches replace the optimizer. They achieve competitive tour lengths but provide no
feasibility guarantee, which is disqualifying when a violated capacity constraint means a vehicle
that physically cannot carry its load.

Bengio, Lodi and Prouvost survey the alternative framing, in which machine learning assists rather
than replaces exact methods --- learning to guide branching, to select heuristics, or to produce
warm starts \cite{bengio2021}. AdaptiveRoute adopts that hybrid stance and applies it at a different
level: the learned component is a fine-tuned language model operating on a natural-language
interface, and the exact solver retains final authority over every plan. The contribution is not a
better tour-construction network; it is an architecture in which an unreliable fast component can be
used safely.

\subsection{Transformer Language Models}

The base model is a causal Transformer \cite{vaswani2017}. The attention mechanism suits this task
because route decisions must be conditioned on relationships among customers, vehicles, capacity
constraints, blocked arcs and the previous plan. An instruction-tuned model was preferred because
the task is contract-driven: the model must follow a structured prompt and emit valid JSON.

\subsection{LoRA and QLoRA}

Full fine-tuning of a 7B-parameter model was not practical on the available hardware. The project
used parameter-efficient fine-tuning with LoRA, which freezes the pretrained weights and injects
trainable low-rank matrices into selected layers \cite{lora}. Training also used QLoRA-style 4-bit
loading, keeping the base model quantized with NF4 and double quantization while adapter weights are
updated \cite{qlora}. The base \path{Qwen/Qwen2.5-7B-Instruct} model remained frozen throughout.

\section{System Architecture}

\subsection{Component Overview}

\begin{verbatim}
Admin / Driver browser
        |
        v
React frontend  (admin console, driver workspace, Leaflet maps)
        |
        v
FastAPI backend
        |
        +--> scenario, driver, route and planning APIs
        +--> agentic routing API
        +--> conversation memory API
        +--> RAG API
        +--> map geometry API
        |
        +--> MongoDB          conversations, messages, context windows,
        |                     agent runs, drivers, routes, planning jobs
        +--> PostgreSQL       document chunks and embeddings (pgvector)
        +--> Pyomo + HiGHS    deterministic CVRP solve, fallback authority
        +--> OSRM             road distance matrix and route geometry
        +--> OpenAI-compatible LLM API
\end{verbatim}

Three interfaces are declared as Python protocols and are swappable at runtime by environment
variable: the routing engine, the route-candidate generator and the event extractor. The repository
layer provides in-memory implementations, so the full test suite runs without a GPU, a database or a
model server.

\subsection{Model Responsibilities}
\label{sec:models}

The system runs two language models with disjoint responsibilities. This separation is deliberate
and is one of the central design decisions.

\begin{table}[h]
\centering
\begin{tabular}{lll}
\toprule
Capability & Model or engine & Role \\
\midrule
Exact planning & Pyomo + HiGHS & Prescriptive source of truth. \\
Interaction and orchestration & Qwen2.5-Coder-32B & Event extraction, route Q\&A, explanation. \\
Route candidate & Qwen2.5-7B + LoRA v5 & Fast approximate replan under disruption. \\
Guardrails & Deterministic validator & Feasibility authority. \\
\bottomrule
\end{tabular}
\end{table}

The rationale is right-sizing. Interpreting an ambiguous driver message is an open-ended language
task that benefits from breadth, so it is handled by a large general model used off the shelf, at
temperature 0.1. Emitting a valid route sequence is a narrow, structured task that benefits from
precision and latency, so it is handled by a small fine-tuned specialist at temperature 0. The
specialist sits on the hot path, which means 32B inference is not paid for every candidate.

The two are configured independently through separate base URLs, model names and temperatures
(\path{ADAPTIVEROUTE_ORCHESTRATOR_*} and \path{ADAPTIVEROUTE_ROUTING_POLICY_*}). The 32B model never
produces a route; the 7B model never addresses a driver.

\subsection{Storage and Deployment}

MongoDB stores document-shaped operational state: conversations, messages, rolling context windows,
agent traces, drivers, operational routes and planning jobs. PostgreSQL with pgvector stores
retrieval chunks and embeddings for documentation-based RAG.

The Docker stack contains the FastAPI backend, the React frontend, MongoDB, PostgreSQL, an optional
OSRM service for road distances and geometry \cite{osrm}, and an optional GPU profile that loads the
LoRA adapter directly in the API process. The OpenAI-compatible LLM API runs outside the Compose
stack and is reached through \path{host.docker.internal}.

\section{Prescriptive Optimization Component}

The Pyomo + HiGHS engine has three roles: it generates supervised labels for training, it provides
the fallback plan when a candidate cannot be accepted, and it is the reference against which the
learned policy is measured.

Let $N$ be the node set, $C \subset N$ the customers, $K$ the vehicles, $d_{ij}$ the travel distance
on arc $(i,j)$, $q_j$ the demand of customer $j$, $Q_k$ the capacity of vehicle $k$, and $B$ the set
of blocked arcs. Let $x_{ijk} \in \{0,1\}$ indicate that vehicle $k$ traverses arc $(i,j)$, and
$y_{jk} \in \{0,1\}$ that vehicle $k$ serves customer $j$.

The objective minimizes total distance:

\[
\min \sum_{k \in K} \sum_{i \in N} \sum_{j \in N, j \neq i} d_{ij}\, x_{ijk}
\]

Each customer is served exactly once:

\[
\sum_{k \in K} y_{jk} = 1
\quad \forall j \in C
\]

Assignment is linked to routing through degree constraints, so that a vehicle enters and leaves a
customer if and only if it serves that customer:

\[
\sum_{i \in N, i \neq j} x_{ijk} = y_{jk},
\qquad
\sum_{i \in N, i \neq j} x_{jik} = y_{jk}
\quad \forall j \in C, \forall k \in K
\]

Capacity is then enforced on the assignment variables:

\[
\sum_{j \in C} q_j\, y_{jk} \leq Q_k
\quad \forall k \in K
\]

Blocked arcs are excluded:

\[
x_{ijk} = 0
\quad \forall (i,j) \in B, \forall k \in K
\]

Subtours are eliminated with MTZ ordering variables $u_{jk}$:

\[
u_{ik} - u_{jk} + |C|\, x_{ijk} \leq |C| - 1
\quad \forall i,j \in C,\ i \neq j,\ \forall k \in K
\]

The implementation adds depot start and end rules and vehicle-usage linking. The linking constraints
above matter: without them, capacity would be decoupled from the routing variables and could be
satisfied vacuously.

\section{Agentic Replanning Workflow}

\subsection{Graph}

The workflow is implemented with LangGraph. Each node is a specialized agent class inheriting from
\path{RoutingWorkflowAgent}, which keeps the graph wiring thin and gives every step a shared contract
for trace and error handling.

\begin{verbatim}
extract_event
      |
      v
solve_base
      |
      v
apply_event
      |
      v
generate_candidate
      |
      v
validate_candidate
      |
      +-- valid ----> compose_response
      |
      +-- invalid --> repair_candidate
                            |
                            +-- valid ----> compose_response
                            |
                            +-- invalid --> solver_fallback
                                                  |
                                                  v
                                          compose_response
\end{verbatim}

If event extraction finds no new event but a context window exists, the workflow answers as a
conversational follow-up rather than forcing a replanning action.

\subsection{Agent Responsibilities}

\begin{longtable}{>{\raggedright\arraybackslash}p{4.2cm}>{\raggedright\arraybackslash}p{10.5cm}}
\toprule
Agent & Responsibility \\
\midrule
\endhead
\path{ExtractEventAgent} & Extract the operational event from user text; report method and confidence. \\
\path{SolveBaseAgent} & Solve the current scenario before mutation; validate the base plan. \\
\path{ApplyEventAgent} & Apply the disruption to the scenario and record the mutation diff. \\
\path{GenerateCandidateAgent} & Generate a route candidate from the solver, an API model or a local LoRA. \\
\path{ValidateCandidateAgent} & Check candidate feasibility against the deterministic contract. \\
\path{RepairCandidateAgent} & Attempt conservative structural repair of an invalid candidate. \\
\path{SolverFallbackAgent} & Re-solve with Pyomo + HiGHS when candidate and repair both fail. \\
\path{ComposeResponseAgent} & Build the response payload with plan, validation, comparison and trace. \\
\bottomrule
\end{longtable}

\subsection{Event Extraction}

Supported events are blocked arc, customer unavailable and priority change. Extraction runs in two
modes: a deterministic rule-based extractor, and an LLM extractor backed by the orchestrator model
with rule-based fallback.

The LLM extractor is not trusted blindly. Its output is parsed into an explicit event object and
rejected if it references nodes that do not exist in the scenario or an unsupported event type. If
the provider is unavailable or returns invalid JSON, extraction falls back to the rule-based path.

\subsection{Candidate Generation Modes}

\begin{table}[h]
\centering
\begin{tabular}{ll}
\toprule
Backend & Behaviour \\
\midrule
\path{solver} & Pyomo + HiGHS as candidate generator. Deterministic; used in tests. \\
\path{api} & The trained LoRA served through an OpenAI-compatible API. \\
\path{local} & The LoRA adapter loaded directly in the API process. \\
\bottomrule
\end{tabular}
\end{table}

Both model-backed modes first check whether the existing plan survives the event. If it does, the
plan is reused and the model is not called at all --- many operational events do not invalidate the
current route.

When a candidate fails validation, the specific violations are serialized back to the model as a
structured correction prompt, for up to three attempts, before the workflow falls through to the
solver.

\section{The Validation Contract}

\subsection{Rules}

Every candidate plan, regardless of origin, is checked against the same deterministic rules:

\begin{itemize}
  \item every active required customer is served exactly once;
  \item inactive or unavailable customers are not served;
  \item no unknown customer or vehicle identifiers appear;
  \item no blocked arc is traversed;
  \item vehicle capacity is respected;
  \item every route starts and ends at the depot.
\end{itemize}

Two properties make this more than a checklist. First, the validator recomputes load and distance
from the scenario rather than trusting the values in the candidate; the model supplies only the
visit sequence, so a hallucinated number cannot reach the output. Second, validation failure is
treated as a prompt rather than an error, feeding the violations back for correction.

\subsection{Repair and Fallback}

Repair is deliberately conservative. It performs only structural cleanups that cannot invent new
optimization decisions: scenario identifier alignment, removal of inactive or duplicated stops,
depot boundary normalization, and recomputation of load and distance. Capacity violations, missing
customers and any genuine rerouting remain the solver's responsibility.

The repaired plan is re-validated before acceptance. If it still fails, the solver fallback produces
an authoritative plan, which is itself validated before being returned.

\subsection{Traceability}

Every node appends trace metadata, exposed in the frontend as an agent execution panel: extraction
confidence and method, candidate source, validation status, repair and fallback decisions, and the
final route source. This is what makes the central claim inspectable rather than asserted --- an
operator can see which component produced the plan they are looking at.

\section{AI Engineering Components}

Beyond the routing policy and the solver, the system carries three components that make the
conversational surface work: a rolling conversation memory, a retrieval layer over project
documentation, and a route question-answering pipeline. They are described here because they
constitute a meaningful part of the implementation and are visible in the operator interface.

\subsection{Conversation Memory and Context Window}

Conversations, messages and agent runs are persisted in MongoDB. On top of the raw message log, the
system maintains a \emph{context window} per conversation --- a compact, bounded summary of state
that survives across turns without replaying the full history into the prompt.

\begin{center}
\begin{tabular}{ll}
\toprule
Field & Content \\
\midrule
\path{summary} & Rolling narrative: latest event, validation, replanning impact. \\
\path{recent_message_ids} & Identifiers of the last $N$ messages. \\
\path{facts} & Durable facts accumulated for the conversation. \\
\path{open_constraints} & Active operational constraints, capped at the last 20. \\
\path{last_event} & Most recent structured event. \\
\path{last_plan} & Most recent accepted plan. \\
\bottomrule
\end{tabular}
\end{center}

The window is rebuilt after each agentic run. The previous summary is carried forward and extended
with the new event, the validation outcome, and the replanning impact --- distance delta and removed
customers --- then truncated to a character budget, keeping the most recent text. Open constraints
accumulate without duplication so that an active blockage remains in scope across turns.

Two details matter operationally. First, when a turn is answered as a conversational follow-up rather
than a replan, the window does not re-record an event, which prevents a question about an existing
disruption from being mistaken for a new one. Second, because the window carries \path{last_plan},
the system can answer questions about the current route without re-solving.

\subsection{Retrieval-Augmented Generation}

The retrieval layer indexes project documentation so that questions about system behaviour can be
answered from written sources rather than model recall.

\paragraph{Ingestion and chunking.}
Ingestion accepts \path{.md}, \path{.txt}, \path{.tex}, \path{.py}, \path{.json}, \path{.yaml},
\path{.yml} and \path{.toml}, skipping build and environment directories. Text is chunked at 1200
characters with 160 characters of overlap. Chunk boundaries snap to the nearest newline, or failing
that the nearest space, but only when that boundary falls past the midpoint of the window --- this
avoids degenerate short chunks while keeping splits off mid-word.

\paragraph{Embeddings.}
Two interchangeable backends implement a common \path{EmbeddingClient} protocol:

\begin{itemize}
  \item \path{hash} (default): a deterministic local embedding. Each token is hashed with BLAKE2b
        into a bucket of a 384-dimensional vector with a signed contribution, and the result is
        L2-normalized. It requires no model server, which makes the retrieval pipeline testable on
        CPU and reproducible across runs.
  \item \path{api}: an OpenAI-compatible \path{/embeddings} endpoint, with returned vectors fitted to
        the configured dimension by truncation or zero-padding.
\end{itemize}

\paragraph{Storage and search.}
A repository protocol has two implementations. The in-memory version scores chunks by cosine
similarity in Python and is used for tests. The production version stores chunks in PostgreSQL with
the \path{pgvector} extension, creating the extension and schema on first use, with the embedding
column typed \path{vector(d)} at the configured dimension. Retrieval uses the cosine distance
operator:

\begin{verbatim}
SELECT ..., 1 - (c.embedding <=> %s::vector) AS score
FROM chunks c
ORDER BY c.embedding <=> %s::vector
LIMIT %s
\end{verbatim}

The subtraction converts distance into a similarity score for ranking and display.

\subsection{Route Question Answering}

Not every driver message is a disruption. Questions such as ``what is my next stop?'' or ``why did
the plan change?'' should be answered without mutating the scenario. These are handled by a separate
pipeline in the conversation service, whose stages are surfaced individually in the agent execution
panel:

\begin{enumerate}
  \item \path{route_lookup} --- resolve the route identifier to a stored route, driver and scenario;
  \item \path{route_fact_builder} --- assemble deterministic facts from the stored plan: stop
        sequence, distances, loads, status;
  \item \path{route_qa_rag} --- retrieve supporting documentation chunks when retrieval is enabled;
  \item \path{route_qa_model} --- compose the natural-language answer with the orchestrator model.
\end{enumerate}

The ordering is deliberate. Route facts are computed from persisted data before the model is called,
so the answer is grounded in the actual plan rather than generated from the conversation alone. If
the model call fails, the pipeline records the failure in the trace and degrades to a fact-based
answer rather than returning nothing.

The classification between question answering and replanning matters: a message implying a route
change must be routed into the replanning workflow, where the validation contract applies. The Q\&A
path is explicitly not an authority on feasibility.

\section{Dataset Generation}

\subsection{Methodology}

Datasets were generated synthetically for scale, repeatability and reliable labels:

\begin{enumerate}
  \item generate a CVRP instance with depot, customers, demands, vehicles, capacities and distance
        matrix;
  \item solve the base plan with Pyomo + HiGHS;
  \item apply an operational event;
  \item re-solve the disrupted scenario;
  \item validate the replanned route;
  \item serialize scenario, event and target into compact SFT JSONL;
  \item convert compact rows into chat-message JSONL for instruction fine-tuning.
\end{enumerate}

The solver therefore labels its own training set. The model learns to imitate prescriptive decisions
while the same engine remains available as a deterministic oracle at runtime.

\subsection{Generated Datasets}

\begin{longtable}{>{\raggedright\arraybackslash}p{5.4cm}rrrr>{\raggedright\arraybackslash}p{3.5cm}}
\toprule
Dataset & Total & Train & Validation & Test & Purpose \\
\midrule
\endhead
\path{training_compact_30k} & 30000 & 24000 & 3000 & 3000 & Initial compact curriculum. \\
\path{training_capacity_tight_50k_parallel} & 50000 & 40000 & 5000 & 5000 & Emphasis on capacity pressure. \\
\path{training_blocked_arc_20k_parallel} & 20000 & 16000 & 2000 & 2000 & Emphasis on blocked arcs. \\
\path{training_curriculum_100k} & 100000 & 80000 & 10000 & 10000 & Consolidated curriculum dataset. \\
\path{training_error_driven_20k} & 20000 & 16000 & 2000 & 2000 & Continuation focused on observed errors. \\
\path{training_error_driven_10k_v6} & 10000 & 8000 & 1000 & 1000 & Additional continuation after v5. \\
\bottomrule
\end{longtable}

The split strategy was 80\% training, 10\% validation, 10\% test. Audits found zero invalid rows in
the principal training datasets.

\section{Fine-Tuning Strategy}

\subsection{Training Algorithm}

Supervised fine-tuning with LoRA/QLoRA, using Hugging Face Transformers for model loading, TRL
\path{SFTTrainer} for supervised fine-tuning, PEFT for adapter management and bitsandbytes for 4-bit
loading. The optimization target is causal language-model loss over the assistant response tokens.

\subsection{Final LoRA Configuration}

\begin{table}[h]
\centering
\begin{tabular}{ll}
\toprule
Parameter & Value \\
\midrule
Base model & \path{Qwen/Qwen2.5-7B-Instruct} \\
Method & LoRA / QLoRA-style loading \\
PEFT type & \path{LORA} \\
Task type & \path{CAUSAL_LM} \\
Rank $r$ & 16 \\
\path{lora_alpha} & 32 \\
\path{lora_dropout} & 0.05 \\
Bias & \path{none} \\
Target modules & \path{q_proj,k_proj,v_proj,o_proj,gate_proj,up_proj,down_proj} \\
Quantization & 4-bit NF4 with double quantization \\
Compute dtype & bfloat16 \\
\bottomrule
\end{tabular}
\end{table}

\section{Training Methodology and Parameters}

Training was conducted in stages, each continuation resuming the previous adapter and evaluated
against the same fixed benchmark before deciding whether to continue or stop.

\begin{longtable}{>{\raggedright\arraybackslash}p{1.1cm}rrrrrr>{\raggedright\arraybackslash}p{3.2cm}}
\toprule
Version & Train rows & Eval rows & Steps & Batch & LR & Time & Notes \\
\midrule
\endhead
v3 & 24000 & 500 & 3000 & 1 $\times$ 8 accum. & $5\cdot10^{-5}$ & 5h53m & First stable full training on the 30K compact dataset. \\
v4 & 56000 & 1000 & 2300 & 8 & $1.7\cdot10^{-5}$ & 3h44m & Hard-data continuation using SDPA and checkpoint resume. \\
v5 & 16000 & 1000 & 1500 & 8 & $1\cdot10^{-5}$ & 2h22m & Error-driven continuation. Selected version. \\
v6 & 8000 & 500 & 800 & 8 & $7\cdot10^{-6}$ & 1h18m & Additional continuation; not selected. \\
\bottomrule
\end{longtable}

Total training time across the four runs was approximately 13 hours on a single RTX 3090. The
decreasing learning rate across stages reflects a deliberate curriculum: broad fitting first, then
progressively finer correction against observed failure modes.

The stabilized setup used \path{--max-seq-length 2048}, \path{--bf16},
\path{--bnb-compute-dtype bfloat16}, \path{--optim adamw_torch},
\path{--lr-scheduler-type cosine}, \path{--max-grad-norm 0.3}, \path{--no-eval-during-train},
\path{dataset_num_proc=1} and \path{dataloader_num_workers=0}. One epoch per run.

The representative v5 command was:

\begin{verbatim}
uv run python scripts/train_lora.py \
  --model-id Qwen/Qwen2.5-7B-Instruct \
  --adapter-path outputs/models/adaptiveroute-qwen2_5-7b-lora-hard70k-v4-b8-sdpa-resume700 \
  --dataset-dir data/chat_training_error_driven_20k \
  --output-dir outputs/models/adaptiveroute-qwen2_5-7b-lora-error20k-v5 \
  --max-seq-length 2048 \
  --epochs 1 \
  --max-steps 1500 \
  --learning-rate 1e-5 \
  --warmup-steps 50 \
  --max-grad-norm 0.3 \
  --lr-scheduler-type cosine \
  --per-device-batch-size 8 \
  --gradient-accumulation-steps 1 \
  --train-limit 16000 \
  --eval-limit 1000 \
  --bf16 \
  --bnb-compute-dtype bfloat16 \
  --optim adamw_torch \
  --save-steps 100 \
  --logging-steps 10 \
  --no-eval-during-train
\end{verbatim}

\section{Evaluation}

\subsection{Training Metrics}

Loss converged to the 0.26--0.28 range after the initial learning stage, and mean token accuracy
stabilized around 0.89. These metrics indicate that the model learned the response format and common
routing patterns, but they are not evidence of operational correctness: the curve is essentially flat
across v4, v5 and v6 while the operational metrics still differ. This is precisely why loss was not
used as the selection criterion.

\begin{table}[h]
\centering
\begin{tabular}{lrrrr}
\toprule
Version & Train loss & Final token accuracy & Tokens processed & Runtime \\
\midrule
v3 & 0.2796 & 0.8922 & $1.614\cdot10^7$ & 21190s \\
v4 & 0.2659 & 0.8915 & $1.241\cdot10^7$ & 13470s \\
v5 & 0.2654 & 0.8911 & $8.081\cdot10^6$ & 8510s \\
v6 & 0.2653 & 0.8917 & $4.307\cdot10^6$ & 4700s \\
\bottomrule
\end{tabular}
\end{table}

\subsection{Operational Evaluation}

All versions were evaluated on the same 1,000-example sample from
\path{data/training_curriculum_100k/sft_test.jsonl}. The primary metric is the feasible-candidate
rate. Exact match was tracked but treated as secondary, because a CVRP instance commonly admits
several equally optimal plans, so matching the reference sequence measures conformity rather than
quality.

\begin{table}[h]
\centering
\begin{tabular}{lrrrrr}
\toprule
Model & Valid JSON & Feasible & Exact match & Capacity violations & Blocked-arc violations \\
\midrule
v3 & 100.0\% & 91.7\% & 4.6\% & 72 & 11 \\
v4 & 100.0\% & 92.9\% & 5.8\% & 61 & 10 \\
v5 & 100.0\% & 94.4\% & 5.6\% & 46 & 9 \\
v6 & 100.0\% & 94.0\% & 5.4\% & 50 & 10 \\
\bottomrule
\end{tabular}
\end{table}

The improvement from v3 to v5 is substantial: feasibility rises 2.7 points and capacity violations
fall by 36\%. The comparison between v5 and v6 requires more care. A difference of 0.4 percentage
points on $n = 1000$ corresponds to four examples; the 95\% confidence interval for a proportion of
0.944 at that sample size is approximately $\pm 1.4$ points, so v5 and v6 are statistically
indistinguishable on feasibility. The honest statement is that the additional continuation produced
no measurable improvement, and v5 was retained as the stopping point rather than v6 being shown to
regress.

v5 was selected because it holds the best observed feasibility, the lowest count in both violation
categories, and because the curve had flattened --- continuing to train without a new measured error
target was not justified.

\subsection{Capacity Boundary}
\label{sec:capacity}

Feasibility on the benchmark describes performance at the size the model was trained for. A separate
benchmark measured how performance degrades as scenario cardinality grows, using strict validation
and a viability threshold of 90\%.

\begin{table}[h]
\centering
\begin{tabular}{rrrl}
\toprule
Customers & Strict valid rate & Avg. candidate latency & Dominant failure mode \\
\midrule
8  & 100\% & 1.59s & --- \\
9  & 100\% & 1.71s & --- \\
10 & 60\%  & 3.99s & customer omitted \\
12 & 33\%  & --- & customer omitted \\
14 & 0\%   & --- & customer omitted \\
\bottomrule
\end{tabular}
\end{table}

The usable boundary is 8--9 customers. Above it, the model omits customers rather than routing them
poorly, and the omissions concentrate on identifiers at or beyond the edge of the training
distribution, such as \path{C10} and above. This is a cardinality generalization limit rather than a
routing-quality problem, and it will not be fixed by prompt engineering; it requires training data at
larger cardinalities.

Two caveats apply. The samples are small --- three to five scenarios per size --- so the intermediate
rates carry wide uncertainty, though the trend and the qualitative failure mode are consistent.
Second, exceeding the boundary does not produce incorrect plans in production: the cascade falls
through to the solver on every request, trading latency for the same guarantee.

\section{Engineering Practices}

The system comprises roughly 11,800 lines of Python across 73 modules, with 24 test files. All
domain and record types are immutable frozen dataclasses. Services and repositories are separated,
with in-memory repository implementations as the default, so the entire test suite runs without a
GPU, a database or a model server.

Continuous integration is split into two jobs: fast API and agentic tests, and solver authority tests
executed in forked subprocesses. The split is necessary because the Pyomo/HiGHS stack can trigger a
native segmentation fault after repeated solver invocations within one Python process.

Driver credentials are stored as bcrypt hashes and authentication issues signed JWTs; CORS origins
are configurable rather than open; requests emit structured JSON logs with correlation identifiers.
A placeholder JWT secret is detected at startup and replaced with a process-local random value, so a
deployment that omits configuration cannot sign tokens with a publicly known key.

\section{Limitations}
\label{sec:limitations}

\subsection{Model Limitations}

\begin{itemize}
  \item The model does not guarantee feasibility or optimality. On the 1,000-example benchmark, v5
        still produced 46 capacity violations and 9 blocked-arc violations.
  \item Exact match is low, though this is partly expected for CVRP.
  \item Training data is entirely synthetic. No real traffic, fleet telemetry, driver constraints,
        service times or geospatial road-network data were used for training.
  \item The usable scenario size is 8--9 customers, as measured in Section~\ref{sec:capacity}.
\end{itemize}

\subsection{Architectural Limitations}

\begin{itemize}
  \item The MTZ formulation has a loose linear relaxation; larger fleets would benefit from a
        stronger branch-and-cut formulation or a dedicated routing solver.
  \item Solver calls still execute inside the request lifecycle. Production use requires a worker
        queue so that optimization leaves the request path.
  \item Validation proves that a plan is correct \emph{for the mutated scenario}. It cannot prove
        that the mutation corresponds to what the driver actually said. Event extraction is the one
        stage the trust boundary does not cover, and a misread message produces a perfectly valid
        answer to the wrong question.
  \item The default embedding backend is the deterministic hash embedding. It makes the retrieval
        pipeline reproducible and testable without a model server, but it is a bag-of-tokens
        projection with no semantic representation: retrieval quality under it is lexical at best.
        Meaningful retrieval requires configuring the API embedding backend, and the current
        deployment has not been evaluated for retrieval precision or recall.
  \item List endpoints support offset pagination but not cursor pagination or server-side filtering.
\end{itemize}

\subsection{Infrastructure Limitations}

\begin{itemize}
  \item Long-running training and generation jobs exposed native failures, including
        \path{exit status 139}, double-free errors, and a \path{torch.nn.Module.train} runtime issue.
  \item FlashAttention could not be installed: the local environment used Python 3.12 with a
        Torch/CUDA 13.0 stack for which no compatible prebuilt wheel was available. The stable
        workaround was SDPA, reduced preprocessing parallelism, evaluation disabled during training,
        and \path{adamw_torch}.
\end{itemize}

\section{Refinement Strategy}

In priority order:

\begin{enumerate}
  \item harden event extraction with confidence thresholds and explicit confirmation before mutating
        a scenario, closing the gap identified in Section~\ref{sec:limitations};
  \item extend the cardinality boundary with curriculum data at 10, 12, 16 and 24 customers, and
        evaluate by scenario size rather than a random split;
  \item move optimization off the request path with a worker queue, persisted solver artifacts and
        time limits tuned per scenario class;
  \item instrument the cascade in production, tracking fallback rate and violation modes as the live
        measure of policy quality;
  \item generate new error-driven datasets from actual validator failures, and train continuations
        only when they beat the incumbent on the fixed benchmark;
  \item evaluate best-of-$N$ candidate generation with deterministic selection;
  \item expand scenario richness: time windows, multiple depots, asymmetric costs, more disruption
        types;
  \item validate against public VRP benchmarks and, ultimately, historical operational routes.
\end{enumerate}

\section{Runtime Contract}

The selected adapter is \path{adaptiveroute-qwen2_5-7b-lora-error20k-v5}, resolved locally under
\path{outputs/models/} with the stable operational alias
\path{outputs/models/adaptiveroute-routing-policy-current} and published as a repository release
asset. Deployments should reference the alias rather than the versioned directory, so that promoting
a future adapter does not require a configuration change.

\begin{verbatim}
driver message
-> event extraction        (orchestrator, with deterministic fallback)
-> base solve              (Pyomo + HiGHS)
-> scenario mutation
-> candidate generation    (LoRA v5, ~1.7s)
-> deterministic validation
-> conservative repair     if invalid
-> Pyomo + HiGHS fallback  if still invalid
-> validated response
\end{verbatim}

The accepted response carries the candidate source, the validation result, any rejected violations,
route distance and vehicle loads, and the final route JSON.

\section{Reproducibility}

The results in this report are reproducible from the repository. This section states the entry
points; operational detail lives in \path{README.md} and \path{docs/DEPLOYMENT.md}.

\paragraph{Obtaining the trained adapter.}
Model artifacts are excluded from version control. The selected LoRA adapter is distributed as a
release asset of the project repository:

\begin{verbatim}
https://github.com/gcarvalho-DataAI/AdaptiveRoute/releases
\end{verbatim}

The asset contains only the adapter weights and their tokenizer configuration, which is a few
hundred megabytes rather than the size of a full checkpoint --- this is the practical benefit of
parameter-efficient fine-tuning. It is not runnable on its own: the frozen base model
\path{Qwen/Qwen2.5-7B-Instruct} is fetched separately from Hugging Face and the adapter is applied on
top of it at load time. Extract the asset under \path{outputs/models/} so that the paths used
throughout this report resolve, or point
\path{ADAPTIVEROUTE_ROUTING_POLICY_LOCAL_ADAPTER_PATH} at wherever it was placed.

Every path in this report can be substituted: the system runs without the adapter by setting
\path{ADAPTIVEROUTE_ROUTING_POLICY_BACKEND=solver}, which uses Pyomo + HiGHS as the candidate
generator. That configuration exercises the full workflow --- extraction, mutation, validation,
repair, fallback --- without a GPU, and is what the test suite uses.

\paragraph{Running the system.}
The default Docker profile requires no GPU and no model weights. It starts the API, the frontend,
MongoDB and PostgreSQL:

\begin{verbatim}
cp .env.docker.example .env.docker
./scripts/docker_up.sh
./scripts/smoke_docker_stack.sh
\end{verbatim}

Road distances and road-snapped geometry require the OSRM profile, which is opt-in and needs its
graph prepared once:

\begin{verbatim}
./scripts/prepare_osrm_nyc.sh
COMPOSE_PROFILES=routing ADAPTIVEROUTE_MAP_ROUTER_BACKEND=osrm ./scripts/docker_up.sh
\end{verbatim}

Without it the system falls back to haversine distances and straight-line geometry. The fallback is
silent by design, so the map legend is the reliable indicator of which backend is active.

\paragraph{Tests.}
The full suite runs against in-memory repositories and requires no GPU, database or model server:

\begin{verbatim}
uv run pytest
\end{verbatim}

Continuous integration splits this into a fast API and agentic job, and a solver job executed with
\path{pytest --forked}, because the Pyomo/HiGHS stack can segfault after repeated solver invocations
within one process.

\paragraph{Reproducing the operational evaluation.}
The feasibility figures in Section~11 are produced in two steps, against the fixed test split:

\begin{verbatim}
uv run python scripts/generate_model_predictions.py \
  --model-path outputs/models/adaptiveroute-qwen2_5-7b-lora-error20k-v5 \
  --dataset data/training_curriculum_100k/sft_test.jsonl \
  --out outputs/predictions/v5_test.jsonl \
  --limit 1000

uv run python scripts/evaluate_predictions.py \
  --dataset data/training_curriculum_100k/sft_test.jsonl \
  --predictions outputs/predictions/v5_test.jsonl \
  --details-out outputs/predictions/v5_eval_details.jsonl
\end{verbatim}

Prediction generation took 20--25 minutes per model on the reference hardware.

\paragraph{Reproducing the capacity benchmark.}
The boundary in Section~11.3 is produced by sweeping scenario size against the served policy:

\begin{verbatim}
uv run python scripts/benchmark_routing_policy_capacity.py \
  --customer-grid 8,9,10,12,14 \
  --samples-per-size 5 \
  --profile balanced \
  --event-types block_arc,customer_unavailable \
  --viability-threshold 0.9 \
  --out outputs/evaluations/routing_policy_capacity.json
\end{verbatim}

Vehicle count defaults to $\lceil \text{customers} / 8 \rceil$ and the seed sequence starts at
10{,}000, so runs are repeatable. Datasets are regenerated deterministically from a seed and are
intentionally excluded from version control.

\section{Training and Test Environment}

\begin{table}[h]
\centering
\begin{tabular}{ll}
\toprule
Item & Observed value \\
\midrule
Operating system & Ubuntu 24.04.4 LTS \\
Kernel & Linux 6.8.0-138-generic \\
CPU & Intel Core i9-14900K, 32 logical CPUs \\
GPU & NVIDIA GeForce RTX 3090 \\
VRAM & 24576 MiB \\
NVIDIA driver & 580.173.02 \\
Compute capability & 8.6 \\
Python & 3.12.3 via \path{uv} \\
Torch & 2.13.0, CUDA 13.0 \\
Transformers & 5.15.1 \\
TRL & 1.10.0 \\
PEFT & 0.20.0 \\
bitsandbytes & 0.50.1 \\
datasets & 5.0.1 \\
Pyomo & 6.10.1 \\
HiGHS/highspy & 1.15.1 \\
\bottomrule
\end{tabular}
\end{table}

\section{Conclusion}

AdaptiveRoute demonstrates that an unreliable fast component can be used safely when the
architecture never depends on it being correct. The fine-tuned routing policy returns feasible
replans 94.4\% of the time at the scenario size it was trained for; the remaining 5.6\% is absorbed
by validation, conservative repair and an exact solver fallback, so the failure rate is a latency
cost rather than a correctness risk.

Three results are worth separating. The first is the policy itself, which is a competent but bounded
artifact: useful to 8--9 customers, measured rather than assumed, with a clear path to extension. The
second is the two-model split, which assigns breadth and precision to different components rather
than asking one model to provide both. The third, and the most durable, is the validation contract:
a single deterministic arbiter that every plan must pass regardless of origin, with the trace exposed
so that an operator can see which component produced what.

The remaining open edge is event extraction. Validation proves a plan is correct for the scenario it
was given; it cannot prove the scenario matches what the driver reported. Closing that gap is the
next substantive piece of work, and it is a guardrail problem rather than a modelling one.

\begin{thebibliography}{9}

\bibitem{dantzig1959}
G. B. Dantzig and J. H. Ramser.
\newblock The Truck Dispatching Problem.
\newblock \textit{Management Science}, 6(1):80--91, 1959.
\newblock \url{https://doi.org/10.1287/mnsc.6.1.80}

\bibitem{mtz}
C. E. Miller, A. W. Tucker and R. A. Zemlin.
\newblock Integer Programming Formulation of Traveling Salesman Problems.
\newblock \textit{Journal of the ACM}, 7(4):326--329, 1960.
\newblock \url{https://doi.org/10.1145/321043.321046}

\bibitem{pointer}
O. Vinyals, M. Fortunato and N. Jaitly.
\newblock Pointer Networks.
\newblock \textit{Advances in Neural Information Processing Systems}, 2015.
\newblock \url{https://arxiv.org/abs/1506.03134}

\bibitem{kool2019}
W. Kool, H. van Hoof and M. Welling.
\newblock Attention, Learn to Solve Routing Problems!
\newblock \textit{International Conference on Learning Representations}, 2019.
\newblock \url{https://arxiv.org/abs/1803.08475}

\bibitem{bengio2021}
Y. Bengio, A. Lodi and A. Prouvost.
\newblock Machine Learning for Combinatorial Optimization: A Methodological Tour d'Horizon.
\newblock \textit{European Journal of Operational Research}, 290(2):405--421, 2021.
\newblock \url{https://arxiv.org/abs/1811.06128}

\bibitem{vaswani2017}
A. Vaswani et al.
\newblock Attention Is All You Need.
\newblock arXiv:1706.03762, 2017.
\newblock \url{https://arxiv.org/abs/1706.03762}

\bibitem{lora}
E. J. Hu et al.
\newblock LoRA: Low-Rank Adaptation of Large Language Models.
\newblock arXiv:2106.09685, 2021.
\newblock \url{https://arxiv.org/abs/2106.09685}

\bibitem{qlora}
T. Dettmers, A. Pagnoni, A. Holtzman and L. Zettlemoyer.
\newblock QLoRA: Efficient Finetuning of Quantized LLMs.
\newblock arXiv:2305.14314, 2023.
\newblock \url{https://arxiv.org/abs/2305.14314}

\bibitem{qwen25}
Qwen Team.
\newblock Qwen2.5 Technical Report and Qwen2.5-7B-Instruct model card.
\newblock \url{https://arxiv.org/abs/2412.15115} and \url{https://huggingface.co/Qwen/Qwen2.5-7B-Instruct}

\bibitem{pyomo}
M. L. Bynum et al.
\newblock Pyomo -- Optimization Modeling in Python.
\newblock Third Edition, Springer, 2021.
\newblock \url{https://www.pyomo.org/citing-pyomo}

\bibitem{highs}
HiGHS.
\newblock High-performance open-source software for linear optimization.
\newblock \url{https://highs.dev/}

\bibitem{osrm}
D. Luxen and C. Vetter.
\newblock Real-time Routing with OpenStreetMap Data.
\newblock \textit{ACM SIGSPATIAL GIS}, 513--516, 2011.
\newblock \url{https://doi.org/10.1145/2093973.2094062}

\end{thebibliography}

\end{document}
